% 第 57 章　什么叫对称
\chapter{什么叫对称}

\step{反常识开场}

冬天窗玻璃上的雪花，几乎每一片都有六个瓣。你拿放大镜看，六个枝杈从中心向六个方向伸出去，彼此之间夹角整整齐齐都是 \(60^\circ\)。可是再看夏天的水滴：一滴挂在草叶尖上的水珠是圆的，你把它绕着中心随便转一个角度——\(1^\circ\) 也好，\(17.3^\circ\) 也好——它看起来都没变。

水变成冰，化学成分一个字都没改，还是那些水分子。可液态的水“转多少角度都不变”，固态的雪花却只剩下“转 \(60^\circ\) 的整数倍才不变”。同样是水，凭什么结冰以后，大自然像是拿着圆规量好了角度，只允许六种朝向？

更奇怪的是，这个“六”并不是雪花独有的。蜜蜂造的蜂巢是六边形的格子，玄武岩冷却收缩出的石柱（比如海边的柱状节理）也是六棱柱，铅笔芯里石墨的一层原子也排成六边形网格。完全不同的物质、完全不同的过程，最后都交出了“六”这个答案。

这些现象你随手就能观察到，但要回答“为什么偏偏是六”“什么才叫对称”，我们平时说的“对称就是好看、就是左右一样”完全不够用。

镜子里还藏着另一个你天天见却很少追问的现象。救护车车头的“救护车”三个字是反着印的，这样前车司机从后视镜里读到的是正字。为什么镜子非要把左右翻过来？你对着镜子举起右手，镜中人举起的是他的左手；可是你低下头，镜中人也低下头，上下怎么就不翻转？镜子明明是一块平平的玻璃，它怎么“知道”该翻左右而不翻上下？这个看似抬杠的问题，恰恰是理解对称最好的入口。

这一章我们要把“对称”这个词从审美词典里拿出来，改造成一件物理学家能用的工具。

\step{先猜一猜}

在往下看之前，先把自己的直觉答案写下来，读完本章再对照。

第一题：下面三样东西，哪一样“最对称”？（a）一张画着标准人脸像的纸；（b）一张画着正方形的纸；（c）一张什么都没画的空白纸。大多数人会选人脸，因为我们从小就被教“人的脸左右对称”。先记下你的选择。

第二题：把两面小镜子竖直立在桌上，镜面相对、夹一个角度，中间放一粒纽扣。你猜镜子里会出现几个纽扣的像？如果夹角是 \(90^\circ\)，是几个？夹角缩小到 \(60^\circ\)，像会变多还是变少？

第三题：物理学家说“自然定律是对称的”。猜一猜这句话可能是什么意思？定律又不是一幅画，它怎么“左右一样”？

\step{动手实验}

\begin{experiment}[两面镜子的万花筒]
找两面可以竖起来的小镜子（化妆镜就行，没有的话两片手机黑屏也能勉强代替），让它们背面相对、镜面打开，像一本立起来翻开的书。在两面镜子中间放一粒纽扣或一截粉笔。

慢慢改变两面镜子的夹角，数镜子里纽扣像的个数。你会发现：

夹角 \(180^\circ\)（两面镜子摊平成一面），只有 1 个像；夹角 \(120^\circ\)，出现 2 个像；夹角 \(90^\circ\)，出现 3 个像；夹角 \(60^\circ\)，出现 5 个像。夹角越小，像越多，而且像的个数总是 \(360^\circ\) 除以夹角再减 1。

再把纽扣贴着夹角的角平分线放，你会发现每个像的动作和真纽扣的动作是联动的：你移动纽扣，所有像按固定的方式跟着动，谁也不能掉队。
\end{experiment}

\begin{figure}[htbp]
\centering
\begin{tikzpicture}[scale=1.1]
  % 两面互成 90 度的镜子（从上方俯视）
  \draw[line width=2pt,blue!60!black] (0,0) -- (3.2,0) node[right]{镜面 1};
  \draw[line width=2pt,blue!60!black] (0,0) -- (0,3.2) node[above]{镜面 2};
  % 物体与三个像
  \fill[red!70!black] (1.0,0.6) circle (2.2pt) node[below right=1pt]{物体};
  \fill[gray] (1.0,-0.6) circle (2.2pt) node[right]{像 1};
  \fill[gray] (-1.0,0.6) circle (2.2pt) node[above]{像 2};
  \fill[gray] (-1.0,-0.6) circle (2.2pt) node[below]{像 3};
  % 反射路径示意
  \draw[-{Stealth},dashed,gray] (1.0,0.6) -- (1.0,0);
  \draw[-{Stealth},dashed,gray] (1.0,0.6) -- (0,0.6);
  \draw (0.55,0) arc (0:90:0.55) node[midway,above right=-2pt]{\(90^\circ\)};
\end{tikzpicture}
\caption{俯视两面夹角 \(90^\circ\) 的镜子：一个物体（俯视）产生三个像。像 1 是物体在镜面 1 里的反射，像 2 是镜面 2 里的反射，像 3 是经过两次反射的“像的像”。三次反射正好把物体“复制”成四个彼此关联的副本——对称操作反复使用，生成一整套副本。}
\end{figure}

先别管“像的个数为什么是 \(360^\circ\) 除以夹角减 1”，注意实验里更深的一件事：镜子里那些像，其实是同一个物体被“复制”出来的副本，而复制它们的规则只有一种——镜面反射。一次反射得到一个像，像再被另一面镜子反射得到下一个像。一种操作，反复使用，就能生成一整串彼此关联的副本。这个实验里藏着本章的核心思想：对称讲的是“操作”，不是“长相”。

顺着“操作”这个线索回家找一圈，你会发现对称操作无处不在：自行车的车轮绕轴转任意角度都差不多（辐条把它变成离散的几十重对称）；电风扇的扇叶转 \(120^\circ\) 重合一次；水龙头的旋钮、门把手、汽车的方向盘，都是工程师故意做成旋转对称的——因为转动的零件最好长得“不挑角度”。反过来，钥匙和锁孔故意做得极不对称，只允许唯一一种插法，这正是“对称性越低、辨识度越高”的反面应用。

\step{旧模型遇到麻烦}

我们对“对称”的旧理解，大致是“左右一样”“看着匀称”“让人舒服”。这个理解在日常生活中够用，但拿来做物理，立刻撞上三堵墙。

第一堵墙：它没法比较。正方形和人脸，哪个更对称？按“好看”的标准，人脸也许赢；可正方形有四条边、四个角全都一样，人脸连左右两边都只是大致相似，鼻孔还经常一大一小。用“美”来量对称，就像用“好听”来量温度——没有刻度，就没有科学。

第二堵墙：它说不清“变的”和“不变的”。雪花的六个瓣明明朝六个方向，长得并不“一样”，可我们偏偏觉得它对称。真正相同的不是六个瓣的位置，而是“你把雪花转 \(60^\circ\) 以后，它和原来分不清”这件事。旧定义盯着物体本身的长相，却漏掉了关键：对称说的是做了某件事之后，有什么没变。

第三堵墙：它管不到最重要的东西。物理学家最关心的是定律：万有引力定律在北京和在纽约一样吗？今天和明天一样吗？实验室朝东摆和朝西摆一样吗？“定律好不好看”是个没有意义的问题，但“换了地点、方向、时间，定律变没变”却是可以拿实验回答的问题。旧定义对此完全使不上劲。

我们需要一个新定义：它要有明确的操作步骤，能数出多少，能应用到物体、状态和定律身上。

\step{发明新概念}

把旧定义倒过来想。不要说“这东西对称，因为它左右一样”，而要说：

先指定一个操作——旋转一个角度、平移一段距离、照一次镜子、等上一段时间；然后问：做完这个操作，我关心的那个东西还能不能和原来分辨开？分辨不开，就说它在这个操作下是对称的。

一句话：对称性就是“做了某种改变之后的不变”。这个定义里有两个主角，一个是“变换”（你做的那件事），一个是“不变性”（做完之后没走样的那部分）。物理学家说“某物具有某种对称性”，永远是在说“某物在某个变换下不变”，两个主角缺一不可。离开变换谈对称，就像离开比赛规则谈输赢。

用这个新定义，我们日常接触的对称操作可以分成几类。

第一类是平移：把东西整体挪一段距离。一条笔直的铁轨沿自己的方向平移，每隔一段枕木间距就重合一次；理想的无限大平面，沿任意方向平移任意距离都不变。

第二类是旋转：绕一个点或一根轴转一个角度。圆任意转都不变，正方形只有转 \(90^\circ\) 的整数倍才不变，雪花只有转 \(60^\circ\) 的整数倍才不变。

第三类是反射：照镜子。蝴蝶、人脸、字母 A 沿中线的反射大致不变；字母 R 一照镜子就变了样。

第四类稍微隐蔽：时间平移。不是移动物体，而是移动“做实验的时刻”。同一个单摆实验，今天上午做和明天上午做，只要条件相同，结果相同。钟表的走时、自由落体的快慢，都不因为你多等了二十四小时而改变。这也是一种对称，而且是最深刻的一种。

还有一个重要的区分：连续对称和离散对称。圆的旋转对称是连续的——允许的角度是一个连续的范围，想转多少转多少。正方形和雪花的旋转对称是离散的——只允许某些特定的角度，像楼梯的台阶，一挡一挡的，中间不许停。镜子实验里像的个数只能是整数，也是离散的特征。

同样重要却常被忽略的，是分清对称的“宿主”。同样一个“转 \(60^\circ\) 不变”，可以长在三种不同的东西身上。

第一种宿主是物体：雪花、正六边形螺母、足球表面重复的图案。这是最容易看见的对称，也是日常语言里“对称”一词的本义。

第二种宿主是状态：一杯静止的水，从哪个方向看都一样，它的静止状态是旋转对称的；一根竖直弹簧的自然下垂状态，绕着弹簧轴旋转也不变。状态不是东西，而是“系统此刻的样子”——本书反复追问的第一个问题“系统此刻是什么状态”，现在多了一种刻画方式：这个状态能承受哪些变换？

第三种宿主是定律：牛顿第二定律写在哪个城市都一样，电磁学规律不因为实验室换了朝向而改一个字。定律的对称不挂在任何具体物体上，它说的是游戏规则本身的公平——不偏心任何地点、方向或时刻。三种宿主中，物理学家最在意的是这一种；前两种对称可以被打破、可以只是近似，第三种对称一旦被实验否定，动摇的是整套理论的地基。

\begin{figure}[htbp]
\centering
\begin{tikzpicture}[scale=1.6]
  % 正六边形（雪花骨架）
  \foreach \i in {0,1,2,3,4,5} {
    \coordinate (v\i) at (60*\i:1);
  }
  \draw[thick] (v0) \foreach \i in {1,2,3,4,5} { -- (v\i) } -- cycle;
  \foreach \i in {0,1,2,3,4,5} {
    \draw[gray] (0,0) -- (v\i);
    \fill (v\i) circle (1.2pt);
  }
  % 60 度旋转箭头
  \draw[-{Stealth},thick,red!70!black] (0:1.45) arc (0:60:1.45)
        node[midway,right=2pt]{\(60^\circ\)};
  \draw[dashed,red!70!black,-{Stealth}] (v0) arc (0:60:1);
  \node at (0,-1.7){转 \(60^\circ\)：顶点互相换位置，图形整体不变};
\end{tikzpicture}
\caption{正六边形的离散旋转对称。转 \(60^\circ\) 后每个顶点都落到相邻顶点的位置上，图形与原来分辨不开；转 \(30^\circ\) 则顶点全部落空。圆没有这种“台阶”，任何角度都行。}
\end{figure}

到这里，我们可以回头给第一题一个出人意料的答案：那张空白纸最对称。它沿任意方向平移任意距离、绕任意点旋转任意角度、沿任意轴反射，都不变——它能“承受”的变换最多。人脸只能承受左右反射（还承受得不严格），正方形只能承受四种旋转和四种反射。直觉告诉我们“空白纸上什么都没有，谈不上对称”，新定义却告诉我们：恰恰是“什么都没有”，才让最多的操作对它无可奈何。对称性度量的不是图案有多丰富，而是有多少种改变无法在它身上留下痕迹。

这个答案还有个实用推论：想比较两件东西谁更对称，就把它们各自“承受得住”的变换列出来数个数。正六边形有 12 个（6 种旋转、6 种反射），正方形有 8 个，长方形只剩 4 个，平行四边形只剩 2 个，不规则四边形只剩“什么都不做”这 1 个。变换集合越大，对称性越高——模糊的美感争议，变成了可以数、可以算的清单。

\step{一句话模型}

\begin{oneline}
对称性 = 做完一个指定操作后，你关心的东西（物体、状态或定律）和原来分辨不开；这个操作就是它的一个对称变换。
\end{oneline}

这个模型像一个“找不同”的游戏：先拍下“之前”的照片，做操作，再拍“之后”的照片，两张照片放到一起对比，找不出区别就算对称。但它不像游戏的地方在于：游戏里两张图本来就该有差别，而物理中的对称要求严格分辨不开；它也不像游戏的地方在于：物理里对比的不只是图案，还可以是运动状态、测量结果乃至定律本身。

\step{数学翻译器}

把“做了操作后不变”翻译成数学，只需要两步：先把操作写成一个规则，再写出“不变”这个条件。

以平面上的旋转为例。纸上每一点用坐标 \((x,y)\) 标记，绕原点旋转角 \(\theta\) 就是一个把旧坐标变成新坐标的规则。“图形在旋转下不变”翻译成：把图形里所有点的坐标按规则换成新坐标，得到的点集和原来的点集是同一个集合。

最简单的检查：正方形绕中心转 \(90^\circ\)，顶点 \((a,a)\) 变成 \((-a,a)\)，\((-a,a)\) 变成 \((-a,-a)\)……四个顶点只是换了个名字，集合没变，所以 \(90^\circ\) 旋转是正方形的对称变换。转 \(45^\circ\) 呢？顶点跑到了原来边的中点方向上，点集变了，所以不是。同样用坐标检查圆：圆心在原点、半径 \(r\) 的圆是满足 \(x^2+y^2=r^2\) 的所有点；旋转不改变任何点到原点的距离，所以新坐标仍满足同一个方程——转任意角度都是对称变换。圆的连续对称和正方形的离散对称，就这样被一行方程区分开了。

平移和反射同样可以翻译。沿 \(x\) 方向平移距离 \(d\)，规则是 \(x'=x+d,\ y'=y\)。铁轨上每隔距离 \(L\) 有一根枕木：取 \(d=L\)，每根枕木恰好落到下一根的位置上，图案不变；取 \(d=L/2\)，枕木全部落空——离散平移对称。理想无限大平面上什么标记都没有，任何 \(d\) 都行——连续平移对称。关于 \(y\) 轴的反射，规则是 \(x'=-x,\ y'=y\)。立刻检查特殊情况：反射两次相当于 \((x,y)\to(-x,y)\to(x,y)\)，回到原处，所以反射是自己的逆操作——这一点和旋转很不一样，旋转一般要做若干次才能回到原处。一只关于 \(y\) 轴对称的蝴蝶，反射后图案不变；一个字母 R，反射后变成镜中反字，图案变了。

再看时间平移。自由落体的规律是下落距离 \(s=\frac{1}{2}gt^2\)。把时间整体推迟 \(T\)，相当于换一个计时起点：用 \(t'=t-T\) 重新计时，规律的写法仍是 \(s=\frac{1}{2}g(t')^2\)，\(g\) 的数值不变。方程的形式在“换计时起点”这个操作下不变，这就是“自由落体定律具有时间平移对称性”的数学说法。立刻用特殊情况检查：取 \(T=0\)，什么都不变，等于没做操作，当然对称；取 \(T\) 为整整一天，说的是今天的实验和明天的实验一致，正是实验可重复的意思；如果 \(g\) 自己随时间漂移，那么换计时起点后方程里就得换一个 \(g'\)，形式被破坏，对称也就没了。

对称性在数学里不只是“检查用”的标签，还是能省掉计算的推理工具。举个力学的例子：一个质量均匀分布的薄圆环，圆心处放一粒小珠子，圆环对珠子的引力指向哪里？你可以逐段计算再求和，但也可以一步算都不算：假如合力指向某个方向，那么把整个装置绕着过圆心、垂直环面的轴转一个小角度，圆环纹丝不动（它是旋转对称的），合力箭头却跟着转了——同一个物理情形竟推出两个不同的答案，矛盾。唯一能自洽的答案是合力为零，因为没有方向的箭头才经得起旋转。这就是对称推理的套路：当“答案若有方向，就讲不出理由”时，答案只能没有方向。后文在电磁学、万有引力的章节里，这种“先对称、后计算”的论证会反复登场。

\begin{mathbridge}
旋转操作可以写得更紧凑。平面上绕原点旋转角 \(\theta\)，新旧坐标的关系是
\[
\begin{pmatrix} x' \\ y' \end{pmatrix}
=
\begin{pmatrix} \cos\theta & -\sin\theta \\ \sin\theta & \cos\theta \end{pmatrix}
\begin{pmatrix} x \\ y \end{pmatrix}.
\]
把中间的方阵记作 \(R(\theta)\)。立刻用特殊情况检查它：\(\theta=0\) 时，\(\cos\theta=1,\sin\theta=0\)，矩阵变成单位矩阵，任何点原地不动——转零度当然等于不转；\(\theta=90^\circ\) 时，点 \((1,0)\) 被送到 \((0,1)\)，正是逆时针转四分之一圈；\(\theta=180^\circ\) 时矩阵是负单位矩阵，每个点都穿过原点落到对面，也符合直觉。

两个旋转可以复合：先转 \(\alpha\) 再转 \(\beta\)，结果是 \(R(\beta)R(\alpha)=R(\alpha+\beta)\)——用余弦、正弦的和角公式就能验证。这个式子说明：旋转这种变换可以“连着做”，连着做的结果还是同类变换。反方向的变换也存在：转 \(\theta\) 的逆操作就是转 \(-\theta\)，两者复合等于不转，\(R(\theta)R(-\theta)=R(0)\)。

数学家把满足四条要求的变换集合叫作“群”：变换可以复合（复合完还在集合里）；有一个“什么都不做”的单位变换；每个变换都有逆变换；复合满足结合律。旋转全体构成一个群，正方形的八个对称变换（四种旋转、四种反射）构成一个只有八个成员的离散群，镜子实验里的反射和它的复合也构成群。“群”这个词听起来高深，其实它只回答一个朴素的问题：把允许的操作一条一条列出来，它们组合起来玩得转吗？整数加法、魔方转动、晶体中原子的排列操作，都是这个问题的不同答案。
\end{mathbridge}

\begin{challenge}
连续的对称变换群（比如平面上所有旋转构成的群，记作 \(SO(2)\)；空间里所有旋转构成的群，记作 \(SO(3)\)）成员无穷多而且连续变化，数学上叫李群。研究它们不必一个一个数成员，而是研究“无穷小变换”——转过极小角度的旋转。所有有限旋转都能由无穷小旋转一层层叠出来，这像微积分里用无穷小位移拼出整条曲线。

离散与连续之间有一道深刻的墙：要把平面用全同的正多边形密铺，只有正三角形、正方形、正六边形做得到，正五边形怎么摆都会留下缝。推广到晶体，能够和平移对称共存的旋转对称只有 2、3、4、6 重。这是数学定理，不是经验总结——所以自然界里没有五重对称的普通晶体。然而 1982 年，谢赫特曼在铝锰合金的电子衍射图样里看到了十次重复的图案，与这条定理公然冲突。后来的解释是：他看到的不是普通晶体，而是“准晶体”——原子排列长程有序但不严格周期，定理的前提不成立，结论自然管不住它。谢赫特曼因此获得 2011 年诺贝尔化学奖。这个例子的教训是：数学定理从不犯错，错的是你以为现实世界满足了定理的全部前提。
\end{challenge}

\step{模型的边界}

“操作之后不变”这个定义干净利落，但用起来有三条边界要守住。

第一条边界：对称永远相对于“你关心的东西”而言。一张人脸照片，对“长相”来说反射后大致不变；对“照片上签名的笔迹”来说反射后完全变了——字成了反字。同一张图，关心这个对称，关心那个不对称。所以物理讨论中必须先说清楚：我们对什么做操作，又要求什么不变。雪花的六重对称指它的外形；它内部某个气泡的位置可能根本没有这个对称。

第二条边界：现实世界里的对称几乎都是近似的。真实的雪花六个瓣并不严格相同，人脸左右从不严格镜像，蜜蜂的巢室也有误差。我们说它们对称，是在“忽略小差异”的模型里说的。这不算偷懒，而是物理的常态——模型抓住主要结构，误差留给更精细的描述。但要警惕另一种极端：因为现实中的对称不严格，就宣布“对称只是人类的美学幻觉”。错了。对称的威力恰恰在于，它是关于定律的陈述，而定律层面的对称可以被检验到极高的精度。

第三条边界，也是最容易踩的坑：不要把“看上去的对称”直接当成“定律的对称”。很长一段时间里，物理学家默认照镜子是世界的基本对称——镜像世界里发生的任何过程，在现实中也应该能发生，这叫宇称对称。1956 年，李政道和杨振宁从理论上提出，弱相互作用主宰的过程里这条对称可能根本不成立；随后吴健雄领导的实验组把钴 60 原子核冷却到极低温度、用磁场排好自旋方向，统计它衰变放出的电子飞向哪个方向。结果清清楚楚：电子明显地偏向与原子核自旋相反的一侧。照镜子以后，这个偏好方向会反过来——镜像世界里发生的衰变，在我们的世界里并不存在。实验事实就是：在弱相互作用中，自然定律在反射变换下并不对称。“宇称在弱相互作用中破缺”是被反复验证的实验结论；至于为什么会这样，则属于理论解释的层面，至今仍是粒子物理的深层问题之一。这个例子给本章的模型划出最锋利的边界：对称不是想当然的默认项，每一种对称都是要交实验答卷的假设。

第四条边界关于推理的分寸：对称论证能排除选项，却不能替你算出数值。圆环中心的引力被对称性钉死为零，但环外某一点的引力是多少，对称只告诉你方向一定在环面与轴线的约束之内，具体大小仍要老老实实算。把“对称所以为零”滥用成“对称所以随便”，是初学者常犯的错误；把一种近似对称当成严格定律来用，则是工程师常犯的错误——真实齿轮的齿总有加工误差，所以再对称的机器也会振动，只是振动的强弱被对称性压到了很小。

\step{回头解释开场现象}

现在回头数一遍开场的问题。

为什么水滴转任意角度都不变？因为把水滴捏成形状的那些因素——表面张力在各个方向上一视同仁，水分子之间的吸引不挑方向。用本章的话说：决定液态水行为的规律在旋转下不变，水滴的状态于是继承了旋转对称，而且是连续的——没有任何因素站出来“钦定”某一个方向。

为什么雪花只有六重？因为水结冰时，分子不是随便堆起来的。冰的晶体里，氧原子之间靠氢键连成六边形的环，环套着环铺成六角网格，相邻氧原子的距离约 \(0.28\) 纳米，网格的夹角锁定在 \(60^\circ\) 和 \(120^\circ\)。雪花长得再大，也是在这样一张微观网格的骨架上一层层往外铺，宏观的六瓣、\(60^\circ\) 的夹角，是微观六角结构放大几千万倍后的直接照片。蜂巢与玄武岩柱的六边形来源不同——那是“等大的圆挤在一起最省边、最省料”的几何结果——但同样可以用对称语言描述：它们从“哪个方向都行”出发，最后落成只允许特定角度的离散图案。

对称在这里经历了三级跳：定律有对称（水分子间的相互作用不挑方向），微观结构只剩离散对称（六角晶格），宏观外形继承了离散对称（六瓣雪花）。这提醒我们区分三种不同的对称：物体的对称（雪花的外形）、状态的对称（一滴静止的水在各个方向上看都一样）、定律的对称（任何地点、方向、时刻做实验，规律不变）。三者常常相关，却不能混为一谈——第 60 章会看到，定律对称而状态不对称，是自然制造丰富结构的基本手法。

至于那个“六”，先做一道估算，感受微观与宏观之间隔了多少层。一片典型雪花质量约 \(1\) 毫克。水的摩尔质量是 \(18\) 克每摩尔，所以 \(1\) 毫克水约含 \(5.6\times10^{-5}\) 摩尔水分子，乘以阿伏伽德罗常数 \(6.02\times10^{23}\)，约为 \(3\times10^{19}\) 个分子。一片雪花直径约 \(5\) 毫米，晶格间距约 \(0.3\) 纳米，也就是说从雪花中心到边缘排着大约一千万个分子。一个包含三千万亿亿个分子的物体，整体的旋转对称性却被两、三个分子搭成的一只六角环严格规定——这就是对称性的乘法：微观的变换规则被复制天文数字次，依然不走样。

镜子实验也能用新语言说清楚了。镜面反射是一个变换，做两次就回到原样。两面夹角为 \(\theta\) 的镜子，反射可以反复复合：物体被镜面 1 反射得像 1，像 1 再被镜面 2 反射得像 3，如此接力下去，每复合一次就生成一个新副本。但副本不能无限多——绕交线一圈总共只有 \(360^\circ\) 的空间，当复合出的副本转满一整圈、落回物体自己的位置时，接力就到头了。算一下就知道像的总数是 \(360^\circ/\theta-1\)：\(90^\circ\) 时得 3 个像，\(60^\circ\) 时得 5 个像，与实验数到的一致。一套看似枯燥的“变换复合规则”，就这样预言了你伸手就能数出来的像的个数。

至于镜子“为什么翻左右不翻上下”，谜底是：镜子根本没翻左右。它翻转的只是垂直于镜面的那个方向——把“指向镜子”变成“离开镜子”。你觉得左右翻了，是因为你想象自己绕竖直轴转 \(180^\circ\) 走到镜后的位置去和镜中人比较，而“翻转镜面方向”与“绕竖直轴转身”这两种操作的效果，恰好在左右上相反、在上下上相同。把比较的方式（也是变换）说清楚，悖论就消失了：不是镜子偏爱左右，而是我们悄悄换了一种转身方式去对比。

还有一个更实际的应用。既然晶体是“同一张对称蓝图复制亿万次”的产物，反过来，只要研究一小块晶体如何散射 X 射线，就能读出整张蓝图：X 射线波长与原子间距同量级，晶体对它来说就是一座规则排列的“衍射光栅”，衍射斑点排成的对称图案直接泄露晶格的对称。1953 年，富兰克林拍下的 DNA 的 X 射线衍射照片呈现特征性的交叉图案，沃森和克里克正是据此推断出 DNA 的双螺旋结构。从“对称图案反推结构”这一思路，今天是材料科学、药物设计的日常工具。另一个天天在你手腕上工作的例子是石英表：石英晶体没有中心反演对称，挤压它时正负电荷中心会错开，表面出现电压（压电效应）；反过来加电压它就形变。让一小片切割方向精确的石英振子以它的固有频率振动，再配合电路计数，就是每秒误差远小于一秒的时钟。对称性的“有”与“没有”，在这里都是可交易的工程资源。工程力学上同理：涡轮叶片按旋转对称排布，是为了让任何一片叶片承受的力学环境与别的叶片分辨不开；齿轮的齿形要沿圆周均匀重复，否则转动时会周期性地“顿”一下。对称在这里不是装饰，是均匀受力的保证。

\step{脑洞实验与挑战题}

本章发明了“对称 = 变换下的不变”这个工具，并试着用它检查了物体、状态和定律。接下来的三章会让这件工具显出真正的威力：第 58 章回答“守恒律从哪里来”，把平移、旋转、时间平移三种对称分别接到动量、角动量、能量守恒上；第 59 章（挑战篇）追问“描述中那些纯属人为的选择（比如电势零点）算不算对称”，走进规范对称；第 60 章则面对“定律对称、世界为什么不单调”的问题，讲对称的破缺。本章的任务只有一个：先把“什么才算对称”说死。

\begin{thought}[搬到宇宙尽头去做实验]
假设你把本章开头那套镜子和纽扣装进飞船，飞到一百亿光年外的星系边上重做实验，顺便还做了自由落体和单摆实验。如果所有结果都和地球实验室里的一模一样，你会说“定律在空间平移下不变”；如果不一样，你就发现了对称的边界。

这不是纯粹的幻想。天文学家观测遥远星系的光谱：氢原子发光的那些特征颜色（谱线）在百亿光年外的星系里依然出现，只是整体因宇宙膨胀向红端移动。把膨胀效应扣除之后，谱线之间的间隔比例与地球实验室里测到的相同。这等于说：百亿光年外的氢原子，遵守和地球实验室里氢原子相同的规律。至少在目力所及的宇宙里，定律的空间平移对称性交出了一份漂亮的答卷。
\end{thought}

\begin{thought}[如果定律明天会变]
反过来想象一个极端世界：自由落体的 \(g\) 每天变一点，电荷之间的库仑力随年份缓慢漂移。在那个世界里，你今天测出的“定律”明天就要作废，实验无法重复，工程图纸隔天过期，连“能量守恒”都无处立足——守恒的前提正是规律不随时间改变，第 58 章会把这个联系讲透。我们习以为常的“实验可以重复”，本身就是一条深刻对称性的日常福利：时间平移对称。
\end{thought}

\begin{puzzles}
\item 字母表里找对称：大写字母 A、H、N、S、O 中，哪些有反射对称？哪些有旋转对称（转 \(180^\circ\) 后不变）？哪一个字母的对称变换最多？提示：先把每个字母允许的变换一条条列出来，再数条数；O 近似看成正圆。
\item 正方形有 8 个对称变换（4 种旋转、4 种反射）。先做一次反射，再旋转 \(90^\circ\)，复合起来的效果是什么？它还是那 8 个变换之一吗？提示：在纸片上画一个正方形，标出四个角 1、2、3、4，动手做一遍，看每个角最终去哪；两次反射复合会得到旋转，这是镜子实验里像的个数有限的根源。
\item 有人说：“既然空白纸最对称，那最对称的宇宙就是最无聊的宇宙。”你怎么看？提示：区分“状态的对称”和“定律的对称”——定律高度对称的同时，状态完全可以丰富多彩；第 60 章会系统讨论这件事。
\item 蜜蜂用六边形筑巢，肥皂泡挤在一起时也常出现六边形边界。这两种“六”与雪花的“六”是同一个原因吗？提示：雪花来自微观晶格的对称；泡沫和蜂巢来自“等大圆盘密铺平面”的几何约束——对称图案相同，背后的变换来源可以完全不同。
\end{puzzles}
